% Content for Chapter 1, Section 1.1
شهدت العملية التعليمية عبر العصور تحولات بنيوية استجابةً للمتغيرات التقنية والاجتماعية؛ فبعد أن كان التعليم يعتمد على التلقين المباشر والمناهج الموحدة التي غالباً ما تُغفل الفروق الفردية، أتاحت التكنولوجيا الرقمية آفاقاً جديدة للتعلم المرن. ومع مطلع القرن الحادي والعشرين، برزت أنظمة إدارة التعلم (\en{LMS}) ومنصات مؤتمرات الفيديو، والتي تسارع تبنيها بشكل كبير إبان جائحة \en{COVID-19}.

إلا أن هذه المنصات الحالية بقيت مجرد ناقل غير تفاعلي للصوت والصورة، حيث يمر المحتوى التعليمي عبرها كتيار بيانات عابر لا يترك أثراً منظماً بعد انتهاء الجلسة، مما يضع عبئاً إدراكياً كبيراً على المعلم في متابعة الفهم، ويترك الطالب في حالة من التلقي غير التفاعلي.

يأتي مشروع \en{IntelliLect} ليمثل الجيل التالي من هذه المنصات، عبر دمج البث المباشر بقدرات الذكاء الاصطناعي التوليديي في الوقت الفعلي. لم يعد الذّكاء الصنعيّ هنا مجرد أداة خارجية، بل "مساعد تدريسي رقمي" يعمل في خلفية المحاضرة؛ فيحول الكلام المنطوق لحظياً إلى نصوص وملخصات، ويضمن دقة المعلومات عبر تقنيات التوليد المعزز بالاسترجاع (\en{RAG})، مما يفرغ المعلم لجودة الشرح ويمنح الطالب تجربة تعليمية تفاعلية ثرية.